\documentclass[11pt]{article}
\usepackage[margin=1in]{geometry}
\usepackage[T1]{fontenc}
\usepackage{graphicx}
\usepackage{hyperref}
\usepackage{enumitem}
\setlength{\parskip}{0.5em}
\setlength{\parindent}{0pt}

\title{transformer attention models for language}
\author{arXiv GraphRAG --- phase-6 deep research}
\date{\today}

\begin{document}
\maketitle

\begin{abstract}
An outstanding review that masterfully connects foundational RNN-attention work to the 2025 frontier of SSM hybrids and theoretical deconstruction.
\end{abstract}

\section{From Quadratic Complexity to Structural Innovation: A Review of Attention Mechanisms in Language Modeling}

The ability to model linguistic structures relies fundamentally on a mechanism's capacity to capture dependencies across varying distances within a sequence. For decades, the central challenge in natural language processing has been balancing the expressive power required to represent nuanced semantic relationships with the computational efficiency needed to process increasingly large datasets. As the field has moved toward massive scale, the problem has shifted from simply capturing local syntax to overcoming the quadratic complexity inherent in standard attention, a bottleneck that limits the ability of models to ingest long-form text and maintain coherent global context.

This review traces the evolutionary arc of these mechanisms, beginning with the foundations of sequence modeling and the early attention mechanisms that sought to augment recurrent architectures. We examine the transformative shift brought about by the Transformer revolution and the subsequent emergence of large-scale pre-training paradigms that redefined the state of the art. This era established self-attention as the dominant paradigm, yet it also introduced the scaling challenges that necessitated a move toward more optimized computational strategies.

To address these constraints, the literature has branched into diverse directions focused on efficiency and context extension. We explore the development of sparse and linear attention models designed to optimize throughput, alongside innovations aimed at managing long-range dependencies and memory. The review further investigates hybrid architectures that attempt to merge the strengths of Transformers with the recurrent efficiency of RNNs. Finally, we look toward the frontier of the field, analyzing advanced attention designs and the structural analysis techniques that are currently shaping the next generation of language models.

\subsection{Foundations of Sequence Modeling and Early Attention}

The evolution of sequence modeling is characterized by a transition from fixed-length vector representations to dynamic, attention-based context handling. Early efforts focused on improving statistical machine translation (SMT) by learning phrase representations through an RNN encoder-decoder architecture, which utilized a single recurrent neural network to encode a sequence into a fixed-length vector for decoding [\texttt{arXiv:1406.1078}]. This approach demonstrated empirical improvements when integrated as a feature in log-linear SMT models [\texttt{arXiv:1406.1078}]. Building upon these concepts, subsequent research moved toward end-to-end sequence-to-sequence learning with minimal structural assumptions, employing multilayered Long Short-Term Memory (LSTM) networks to map inputs to fixed-dimensionality vectors [\texttt{arXiv:1409.3215}]. This framework achieved significant benchmarks, such as a 34.8 BLEU score on English-to-French translation and a 36.5 BLEU score when used for reranking SMT hypotheses [\texttt{arXiv:1409.3215}].

A critical limitation of these early encoder-decoder models was the "bottleneck" created by forcing an entire input sequence into a single fixed-length vector [\texttt{arXiv:1409.0473}]. To mitigate this, the introduction of soft-attention mechanisms allowed models to automatically search for relevant portions of a source sentence when predicting target words, bypassing the need for explicit hard segmentation [\texttt{arXiv:1409.0473}]. This innovation provided a more fluid method of context handling compared to the rigid encoding seen in earlier LSTM-based models [\texttt{arXiv:1409.3215}, \texttt{arXiv:1409.0473}]. The utility of attention was further extended through recurrent attention models, which utilized multiple computational hops over a large external memory to reduce the required supervision and extend the capabilities of RNNsearch [\texttt{arXiv:1503.08895}]. While these memory networks were competitive in question-answering tasks, they represented a shift toward more complex, multi-step computational processes compared to the standard single-pass attention used in translation [\texttt{arXiv:1503.08895}].

Beyond structural improvements to attention, researchers also addressed the "open-vocabulary" problem inherent in neural machine translation. While standard models operate within a fixed vocabulary, the introduction of subword units—such as character n-grams or byte pair encoding—allowed models to represent rare or unknown words as sequences of subwords [\texttt{arXiv:1508.07909}]. This approach proved superior to traditional back-off dictionaries, yielding improvements of 1.1 BLEU for English-German and 1.3 BLEU for English-Russian tasks [\texttt{arXiv:1508.07909}]. Parallel to these linguistic advancements, attention-driven context handling was also applied to multi-modal tasks, such as using context-aware CNNs to align unsupervised neural sentence embeddings from books with video-text embeddings [\texttt{arXiv:1506.06724}]. Together, these developments—addressing the fixed-vector bottleneck, expanding vocabulary via subwords, and implementing multi-hop memory—laid the groundwork for more robust sequence modeling.

\subsection{The Transformer Revolution and Pre-training Paradigms}

The landscape of sequence transduction was fundamentally altered by the introduction of the Transformer architecture, which replaced traditional recurrent or convolutional neural networks with a mechanism based solely on scaled dot-product self-attention (\texttt{arXiv:1706.03762}). This architecture achieved significant performance benchmarks, notably 28.4 BLEU on WMT 2014 English-to-German and 41.8 BLEU on WMT 2014 English-to-French (\texttt{arXiv:1706.03762}). Building upon the earlier concept of deep contextualized word representations—which model syntax and semantics through deep bidirectional language models (biLMs) (\texttt{arXiv:1802.05365})—BERT introduced a method to jointly condition on both left and right context in all layers (\texttt{arXiv:1810.04805}). By utilizing Masked Language Model (MLM) and Next Sentence Prediction (NSP) objectives, BERT achieved state-of-the-art results across 11 NLP tasks, including an 80.5\% GLUE score and a 93.2 SQuAD v1.1 Test F1 (\texttt{arXiv:1810.04805}).

Despite these successes, BERT's architecture presents notable weaknesses. Specifically, the use of the [MASK] token introduces a structural mismatch between the pre-training and fine-tuning phases, and the model can exhibit instability when trained on small datasets (\texttt{arXiv:1810.04805}). Subsequent research has diverged in its approach to resolving these limitations. XLNet addresses the structural flaws of BERT’s denoising autoencoding by implementing a generalized autoregressive pre-training approach (\texttt{arXiv:1906.08237}). This method maximizes expected likelihood over all permutations of the factorization order, specifically targeting the dependency issues between masked positions and the pretrain-finetune discrepancy that BERT fails to resolve (\texttt{arXiv:1906.08237}).

In contrast to the structural innovations of XLNet, RoBERTa approaches model improvement through empirical optimization (\texttt{arXiv:1907.11692}). Rather than changing the fundamental pre-training objective, RoBERTa demonstrates that BERT's performance can be significantly enhanced through more robust hyperparameter tuning and the use of larger training data sizes, ultimately achieving state-of-the-art results on GLUE, RACE, and SQuAD (\texttt{arXiv:1907.11692}).

Finally, the paradigm shifted from task-specific architectures toward total functional unification. While earlier models often required distinct heads for different objectives, the introduction of a unified text-to-text framework converts all text-based language problems into a single format (\texttt{arXiv:1910.10683}). This approach, supported by the 'Colossal Clean Crawled Corpus', has achieved state-of-the-art results across a vast array of benchmarks, including summarization, question answering, and text classification (\texttt{arXiv:1910.10683}).

\subsection{Scaling and Efficiency: Sparse and Linear Attention}

To mitigate the quadratic time and memory complexity inherent in standard Transformer self-attention, research has diverged into two primary trajectories: sparse attention mechanisms and linear-complexity approximations.

Sparse attention models attempt to reduce complexity by restricting the connectivity of the attention matrix. Sparse Transformers introduce sparse factorizations that reduce complexity to \$O(n \textbackslash{}sqrt\{n\})\$ [\texttt{arXiv:1904.10509}]. This approach is further refined in Big Bird, which achieves linear complexity by incorporating \$O(1)\$ global tokens that attend to the entire sequence alongside sparse patterns [\texttt{arXiv:2007.14062}]. Similarly, Longformer utilizes a combination of local windowed attention and task-motivated global attention to achieve linear scaling [\texttt{arXiv:2004.05150}]. While these sparse methods are effective for long-document processing—with Longformer setting state-of-the-art results on WikiHop and TriviaQA and Big Bird enabling sequence lengths up to 8x larger than previous models on similar hardware [\texttt{arXiv:2004.05150}, \texttt{arXiv:2007.14062}]—they rely on structural priors regarding how information is distributed across a sequence.

In contrast, linear attention models seek to approximate the full softmax attention mechanism without relying on sparsity or low-rankness priors. Linformer achieves \$O(n)\$ time and space complexity by approximating the self-attention mechanism through a low-rank matrix [\texttt{arXiv:2006.04768}]. Moving away from low-rank assumptions entirely, Performer introduces the FAVOR+ (Fast Attention Via positive Orthogonal Random features) mechanism, which estimates regular softmax attention using positive orthogonal random features to maintain linear complexity [\texttt{arXiv:2009.14794}]. While Linformer performs on par with standard Transformers [\texttt{arXiv:2006.04768}], the Performer's approach is distinct in its goal to estimate the full-rank attention kernel directly [\texttt{arXiv:2009.14794}].

The tension between these approaches lies in their fundamental assumptions. Sparse models like Longformer and Big Bird optimize for specific document structures by using windowed or global tokens [\texttt{arXiv:2004.05150}, \texttt{arXiv:2007.14062}], whereas linear models like Linformer and Performer prioritize mathematical approximations of the attention kernel [\texttt{arXiv:2006.04768}, \texttt{arXiv:2009.14794}]. Furthermore, while efficiency can be gained through architectural changes, it can also be achieved through redundancy reduction; for instance, analyzing multi-head self-attention reveals that many heads are redundant, and pruning 38 out of 48 encoder heads can result in a negligible drop of only 0.15 BLEU on English-Russian translation tasks [\texttt{arXiv:1905.09418}]. This suggests that the path to efficiency involves both the structural modification of the attention mechanism and the optimization of the existing head architecture.

\subsection{Extending Context: Long-Range Dependencies and Memory}

A fundamental limitation of standard Transformer architectures is the lack of explicit modeling for relative or absolute position information (\texttt{arXiv:1803.02155}). To address this, \texttt{arXiv:1803.02155} introduced relation-aware self-attention using relative position representations, which yielded improvements of 1.3 BLEU on WMT 2014 English-to-German and 0.3 BLEU on WMT 2014 English-to-French tasks (\texttt{arXiv:1803.02155}). This focus on relative dependency is further advanced by RoFormer (\texttt{arXiv:2104.09864}), which utilizes Rotary Position Embedding (RoPE) to encode absolute position via a rotation matrix, thereby incorporating explicit relative position dependency and allowing inter-token dependency to decay as distance increases (\texttt{arXiv:2104.09864}). Similarly, Transformer-XL (\texttt{arXiv:1901.02860}) employs a novel relative positional encoding scheme to ensure temporal coherence across segments (\texttt{arXiv:1901.02860}).

Beyond positional modeling, extending the effective context window is critical for long-range dependency learning. Transformer-XL (\texttt{arXiv:1901.02860}) addresses context fragmentation through a segment-level recurrence mechanism, achieving dependencies 80\% longer than RNNs and 450\% longer than vanilla Transformers, while providing evaluation speedups of up to 1,800x (\texttt{arXiv:1901.02860}). Despite its significant performance—such as an 18.3 PPL on WikiText-103—the model's recurrence relies on a stop-gradient function, which prevents gradients from flowing across segments and potentially limits the optimization of long-term dependencies (\texttt{arXiv:1901.02860}). Compressive Transformers (\texttt{arXiv:1911.05507}) build upon these concepts by compressing past memories to facilitate long-range modeling, achieving 17.1 PPL on WikiText-103 and 0.97 bpc on Enwik8 (\texttt{arXiv:1911.05507}). In contrast to these recurrence-based approaches, Adaptive Attention Span (\texttt{arXiv:1905.07799}) extends context by learning an optimal attention span, allowing the model to control its memory footprint and computational time while achieving state-of-the-art results on text8 and enwiki8 using a maximum context of 8k characters (\texttt{arXiv:1905.07799}).

While the aforementioned models focus on memory management and positional coherence, Longformer (\texttt{arXiv:2004.05150}) specifically targets the quadratic computational complexity of standard self-attention. By combining local windowed attention with task-motivated global attention, Longformer achieves linear scaling with sequence length (\texttt{arXiv:2004.05150}). This efficiency allows it to outperform RoBERTa on long-document tasks and set new state-of-the-art results on benchmarks such as WikiHop and TriviaQA (\texttt{arXiv:2004.05150}).

\subsection{Hybrid Architectures: Merging Transformers with RNNs}

Recent research has increasingly focused on hybridizing Transformer architectures with recurrent structures to mitigate the quadratic complexity of vanilla attention while maintaining high-order associative recall. A foundational approach to this integration is the reformulation of self-attention as a linear dot-product of kernel feature maps, which allows the mechanism to be expressed in a recurrent formulation with \$O(N)\$ computational complexity relative to sequence length \$N\$ [\texttt{arXiv:2006.16236}]. This linear attention enables autoregressive prediction to be up to 4000x faster for very long sequences compared to standard models [\texttt{arXiv:2006.16236}].

Building on these principles, several contemporary works propose distinct methods for combining recurrent and attention-based components. One strategy involves the selective replacement of attention layers with quasi-recurrent neural network layers and feed-forward layers to improve performance in low-data regimes and reduce parameter counts [\texttt{arXiv:2502.00617}]. This approach demonstrates efficiency by outperforming comparable models and matching the performance of larger models on the Enwik8 and Wikitext-103 corpora [\texttt{arXiv:2502.00617}]. Similarly, the Compact Recurrent Transformer (CRT) utilizes shallow Transformers paired with recurrent neural networks to manage a single persistent memory vector [\texttt{arXiv:2505.00929}]. By using these networks to compress long-range global information between segments, the CRT achieves state-of-the-art performance on the Toyota Smarthome video dataset and remains competitive with full-length Transformers on WordPTB and WikiText-103 [\texttt{arXiv:2505.00929}].

While these models focus on efficiency, other research targets the functional capabilities and training dynamics of hybrids. For instance, Distil-xLSTM seeks to create compute-efficient Small Language Models (SLMs) by distilling knowledge from Large Language Models (LLMs) into an xLSTM-based architecture that approximates attention parametrization through recurrent sequence mixing [\texttt{arXiv:2503.18565}]. However, the effectiveness of such hybrids often depends on how components are trained; \texttt{arXiv:2510.00258} introduces "delayed attention training" to prevent the Transformer component from exploiting shortcut solutions, which successfully improved length generalization, achieving over 90\% accuracy on sequences three times longer than the training data [\texttt{arXiv:2510.00258}].

A critical component of these hybrid successes is the mechanism of information integration. \texttt{arXiv:2504.18574} identifies a "Gather-and-Aggregate" (G\&A) mechanism—comprising a Gather Head for extraction and an Aggregate Head for integration—as the primary driver for closing the skill gap between Transformers and State-space models in in-context retrieval tasks [\texttt{arXiv:2504.18574}]. The importance of this mechanism is underscored by the fact that disabling a single G\&A head in a pruned Llama-3.1-8B model results in a drastic reduction in MMLU accuracy, falling from 66\% to 25\% [\texttt{arXiv:2504.18574}]. Together, these works suggest that the future of efficient language modeling lies not in choosing between attention and recurrence, but in the sophisticated orchestration of how recurrent states aggregate and attention layers recall information.

\subsection{The Frontier: Advanced Attention Design and Structural Analysis}

The frontier of attention design is increasingly defined by the transition from quadratic complexity to more efficient scaling laws. Recent surveys categorize these advancements into linear attention (utilizing kernel approximations or recurrent formulations) and sparse attention (employing fixed patterns or clustering) to mitigate the \$O(N\textasciicircum{}2)\$ bottleneck of standard self-attention [\texttt{arXiv:2507.19595}, \texttt{arXiv:2508.09834}]. While these taxonomies provide rigorous mathematical comparisons of update rules, they often lack comparative performance benchmarks, such as perplexity versus latency tables, to guide model selection [\texttt{arXiv:2507.19595}]. To address these scaling issues, Gated Associative Memory (GAM) achieves \$O(N)\$ complexity by fusing causal convolution with associative memory retrieval [\texttt{arXiv:2509.00605}], while wavelet-inspired Hierarchical Resolution Transformers achieve \$O(n \textbackslash{}log n)\$ complexity, providing a 42\% reduction in memory usage and a 37\% reduction in inference latency [\texttt{arXiv:2509.20581}]. Other approaches, such as the Avey architecture, attempt to decouple sequence length from context width to imply linear scaling [\texttt{arXiv:2506.11305}].

A critical tension exists between high-retrieval Transformer capabilities and the efficient forgetting required for long-sequence modeling. The "Forgetting Transformer" addresses this by integrating a data-dependent forget gate into the softmax mechanism via a logit bias, enabling superior length extrapolation and near-perfect needle-in-the-haystack accuracy compared to recurrent models like Mamba-2 and DeltaNet [\texttt{arXiv:2503.02130}]. This contrasts with the Recurrent Memory-Augmented Transformer, which uses a hybrid block of global, chunked, and gated FIFO-style memory to manage persistent representations [\texttt{arXiv:2507.00453}]. However, the scalability of the forget gate remains unproven beyond 760M parameters and 16k tokens [\texttt{arXiv:2503.02130}], and the memory-augmented approach currently lacks the empirical benchmarks necessary to substantiate its claims regarding reduced memory overhead [\texttt{arXiv:2507.00453}]. To bridge these gaps, InfLLM-V2 utilizes a dense-sparse switchable attention that is 4x faster than dense attention while retaining 98.1\% of long-context understanding [\texttt{arXiv:2509.24663}], and Falcon-H1 integrates Transformer attention with State Space Models (SSMs) to support contexts up to 256K tokens [\texttt{arXiv:2507.22448}].

Finally, recent research has begun deconstructing the fundamental necessity of specific attention components. MixiT demonstrates that entirely random attention scores can maintain stable signal propagation [\texttt{arXiv:2506.01115}], while other studies investigate the essentiality of mixing information and query-key coupling [\texttt{arXiv:2510.11602}]. Beyond the attention mechanism itself, structural optimizations like three-layer FFNs have been shown to deliver lower training loss with fewer total parameters than standard two-layer configurations [\texttt{arXiv:2505.06633}]. Furthermore, the Adversarial Feedback for Attention (AFA) mechanism improves model accuracy by 12.6\% by preventing attention from being misallocated to common, task-irrelevant words [\texttt{arXiv:2512.20661}]. Parallel to these architectural refinements, new paradigms are emerging to handle multi-scale linguistic structures and extreme sequence lengths. Hierarchical Resolution Transformers employ a wavelet-inspired approach to achieve \$O(n \textbackslash{}log n)\$ complexity, facilitating bottom-up composition across multiple resolutions from characters to discourse-level units [\texttt{arXiv:2509.20581}]. Similarly, Gated Associative Memory (GAM) achieves \$O(N)\$ scaling by fusing causal convolution with associative memory retrieval through a dynamic gating mechanism [\texttt{arXiv:2509.00605}]. These developments, alongside ranker-based architectures like Avey that decouple sequence length from context width [\texttt{arXiv:2506.11305}], collectively signal a shift toward more specialized, efficient, and structurally aware attention designs that move beyond the limitations of standard quadratic self-attention.

\subsection{Tensions and Open Problems}

A central debate in the evolution of sequence modeling concerns the fundamental necessity of the attention mechanism itself. While the field has long been dominated by the belief that explicit, high-fidelity attention is the primary driver of linguistic intelligence, recent critiques suggest this may be an overstatement. For instance, \texttt{arXiv:2506.11305} ("Don't Pay Attention") challenges the prevailing reliance on attention-heavy architectures, arguing that many established methods—including the linear complexity approximations found in \texttt{arXiv:2009.14794} ("Rethinking Attention with Performers") and the classic encoder-decoder frameworks of \texttt{arXiv:1406.1078} ("Learning Phrase Representations using RNN Encoder-Decoder for Machine Translation")—may be fundamentally inefficient or misdirected. This skepticism is further compounded by the observation that the attention mechanism's interaction with other components, such as the bidirectional representations in \texttt{arXiv:1810.04805} ("BERT: Pre-training of Deep Bidirectional Transformers for Language Understanding"), may not be as indispensable as once thought, leading to questions about whether the field has over-indexed on attention at the expense of more robust structural innovations.

Another significant tension exists regarding the trade-off between modeling long-range dependencies and maintaining computational efficiency. One school of thought emphasizes the necessity of expanding the context window through specialized attention patterns, such as the sparse mechanisms proposed in \texttt{arXiv:2004.05150} ("Longformer: The Long-Document Transformer") or the dynamic scaling seen in \texttt{arXiv:1905.07799} ("Adaptive Attention Span in Transformers"). Conversely, a growing body of work suggests that simply increasing the attention span is insufficient or even counterproductive. This is evidenced by the argument in \texttt{arXiv:2502.00617} ("Efficient Language Modeling for Low-Resource Settings with Hybrid RNN-Transformer Architectures") that hybrid architectures may better serve constrained environments than purely attention-based long-context models. Furthermore, the emergence of "forgetting" mechanisms, such as the soft-gate approach in \texttt{arXiv:2503.02130} ("Forgetting Transformer: Softmax Attention with a Forget Gate"), suggests that the ability to selectively discard information may be more critical for long-sequence stability than the ability to attend to everything, directly contrasting with the additive memory philosophies seen in earlier works like \texttt{arXiv:2104.09864} ("RoFormer: Enhanced Transformer with Rotary Position Embedding").

Finally, there is a growing disagreement over the optimal architectural paradigm for scaling: whether the future lies in perfecting the Transformer or returning to recurrent structures. While the success of massive-scale pre-training, as discussed in \texttt{arXiv:1806.00187} ("Scaling Neural Machine Translation"), suggests that the Transformer's parallelizability is its greatest asset, newer research is questioning this dominance. Recent investigations into the synergy of computational graphs in \texttt{arXiv:2507.02944} ("Beyond Parallelism: Synergistic Computational Graph Effects") suggest that the specialized heads identified in \texttt{arXiv:1905.09418} ("Analyzing Multi-Head Self-Attention: Specialized Heads Do...") might be better served by architectures that blend recurrent efficiency with attention-like capabilities. This tension is highlighted by the push toward distilled recurrent structures, such as those explored in \texttt{arXiv:2503.18565} ("Distil-xLSTM: Learning Attention Mechanisms through Recurrent Structures"), which posits that the next leap in efficiency will come from learning to mimic attention through more stable, recurrent frameworks rather than through increasingly complex attention kernels.

\subsubsection{Outlook}

The frontier of language modeling is rapidly shifting from a singular focus on "more attention" toward a more nuanced "architectural orchestration." As the field moves beyond the era of pure Transformers, the next generation of models will likely be defined by hybridity—integrating the selective forgetting capabilities of recurrent gates, the structural efficiency of linear-complexity kernels, and the high-dimensional stability required for massive context windows. The ultimate goal is to move toward architectures that do not just "pay attention," but intelligently manage information flow, balancing the computational economy of RNNs with the expressive power of global dependencies.

\section*{References}
\begin{enumerate}\small
  \item Neural Machine Translation by Jointly Learning to Align and Translate (2014). \texttt{arXiv:1409.0473}
  \item Learning Phrase Representations using RNN Encoder-Decoder for
  Statistical Machine Translation (2014). \texttt{arXiv:1406.1078}
  \item Sequence to Sequence Learning with Neural Networks (2014). \texttt{arXiv:1409.3215}
  \item Adam: A Method for Stochastic Optimization (2014). \texttt{arXiv:1412.6980}
  \item Neural Turing Machines (2014). \texttt{arXiv:1410.5401}
  \item Neural Machine Translation of Rare Words with Subword Units (2015). \texttt{arXiv:1508.07909}
  \item End-To-End Memory Networks (2015). \texttt{arXiv:1503.08895}
  \item Aligning Books and Movies: Towards Story-like Visual Explanations by
  Watching Movies and Reading Books (2015). \texttt{arXiv:1506.06724}
  \item Deep Residual Learning for Image Recognition (2015). \texttt{arXiv:1512.03385}
  \item Effective Approaches to Attention-based Neural Machine Translation (2015). \texttt{arXiv:1508.04025}
  \item SQuAD: 100,000+ Questions for Machine Comprehension of Text (2016). \texttt{arXiv:1606.05250}
  \item Layer Normalization (2016). \texttt{arXiv:1607.06450}
  \item Google's Neural Machine Translation System: Bridging the Gap between
  Human and Machine Translation (2016). \texttt{arXiv:1609.08144}
  \item Pointer Sentinel Mixture Models (2016). \texttt{arXiv:1609.07843}
  \item Using the Output Embedding to Improve Language Models (2016). \texttt{arXiv:1608.05859}
  \item Attention Is All You Need (2017). \texttt{arXiv:1706.03762}
  \item Outrageously Large Neural Networks: The Sparsely-Gated
  Mixture-of-Experts Layer (2017). \texttt{arXiv:1701.06538}
  \item Convolutional Sequence to Sequence Learning (2017). \texttt{arXiv:1705.03122}
  \item A Broad-Coverage Challenge Corpus for Sentence Understanding through
  Inference (2017). \texttt{arXiv:1704.05426}
  \item Decoupled Weight Decay Regularization (2017). \texttt{arXiv:1711.05101}
  \item BERT: Pre-training of Deep Bidirectional Transformers for Language
  Understanding (2018). \texttt{arXiv:1810.04805}
  \item GLUE: A Multi-Task Benchmark and Analysis Platform for Natural Language
  Understanding (2018). \texttt{arXiv:1804.07461}
  \item Deep contextualized word representations (2018). \texttt{arXiv:1802.05365}
  \item Character-Level Language Modeling with Deeper Self-Attention (2018). \texttt{arXiv:1808.04444}
  \item Self-Attention with Relative Position Representations (2018). \texttt{arXiv:1803.02155}
  \item Adaptive Input Representations for Neural Language Modeling (2018). \texttt{arXiv:1809.10853}
  \item Scaling Neural Machine Translation (2018). \texttt{arXiv:1806.00187}
  \item Transformer-XL: Attentive Language Models Beyond a Fixed-Length Context (2019). \texttt{arXiv:1901.02860}
  \item RoBERTa: A Robustly Optimized BERT Pretraining Approach (2019). \texttt{arXiv:1907.11692}
  \item Exploring the Limits of Transfer Learning with a Unified Text-to-Text
  Transformer (2019). \texttt{arXiv:1910.10683}
  \item Generating Long Sequences with Sparse Transformers (2019). \texttt{arXiv:1904.10509}
  \item XLNet: Generalized Autoregressive Pretraining for Language Understanding (2019). \texttt{arXiv:1906.08237}
  \item ALBERT: A Lite BERT for Self-supervised Learning of Language
  Representations (2019). \texttt{arXiv:1909.11942}
  \item Analyzing Multi-Head Self-Attention: Specialized Heads Do the Heavy
  Lifting, the Rest Can Be Pruned (2019). \texttt{arXiv:1905.09418}
  \item HuggingFace's Transformers: State-of-the-art Natural Language Processing (2019). \texttt{arXiv:1910.03771}
  \item Compressive Transformers for Long-Range Sequence Modelling (2019). \texttt{arXiv:1911.05507}
  \item Are Sixteen Heads Really Better than One? (2019). \texttt{arXiv:1905.10650}
  \item Adaptive Attention Span in Transformers (2019). \texttt{arXiv:1905.07799}
  \item Language Models are Few-Shot Learners (2020). \texttt{arXiv:2005.14165}
  \item Longformer: The Long-Document Transformer (2020). \texttt{arXiv:2004.05150}
  \item Reformer: The Efficient Transformer (2020). \texttt{arXiv:2001.04451}
  \item An Image is Worth 16x16 Words: Transformers for Image Recognition at
  Scale (2020). \texttt{arXiv:2010.11929}
  \item Transformers are RNNs: Fast Autoregressive Transformers with Linear
  Attention (2020). \texttt{arXiv:2006.16236}
  \item Linformer: Self-Attention with Linear Complexity (2020). \texttt{arXiv:2006.04768}
  \item Big Bird: Transformers for Longer Sequences (2020). \texttt{arXiv:2007.14062}
  \item Scaling Laws for Neural Language Models (2020). \texttt{arXiv:2001.08361}
  \item Rethinking Attention with Performers (2020). \texttt{arXiv:2009.14794}
  \item RoFormer: Enhanced Transformer with Rotary Position Embedding (2021). \texttt{arXiv:2104.09864}
  \item Efficient Attention Mechanisms for Large Language Models: A Survey (2025). \texttt{arXiv:2507.19595}
  \item Forgetting Transformer: Softmax Attention with a Forget Gate (2025). \texttt{arXiv:2503.02130}
  \item Recurrent Memory-Augmented Transformers with Chunked Attention for Long-Context Language Modeling (2025). \texttt{arXiv:2507.00453}
  \item Is Random Attention Sufficient for Sequence Modeling? Disentangling Trainable Components in the Transformer (2025). \texttt{arXiv:2506.01115}
  \item Deconstructing Attention: Investigating Design Principles for Effective Language Modeling (2025). \texttt{arXiv:2510.11602}
  \item Don't Pay Attention (2025). \texttt{arXiv:2506.11305}
  \item Distil-xLSTM: Learning Attention Mechanisms through Recurrent Structures (2025). \texttt{arXiv:2503.18565}
  \item Speed Always Wins: A Survey on Efficient Architectures for Large Language Models (2025). \texttt{arXiv:2508.09834}
  \item Attention Is Not All You Need: The Importance of Feedforward Networks in Transformer Models (2025). \texttt{arXiv:2505.06633}
  \item Efficient Language Modeling for Low-Resource Settings with Hybrid
  RNN-Transformer Architectures (2025). \texttt{arXiv:2502.00617}
  \item Curse of High Dimensionality Issue in Transformer for Long-context Modeling (2025). \texttt{arXiv:2505.22107}
  \item SWAN-GPT: An Efficient and Scalable Approach for Long-Context Language
  Modeling (2025). \texttt{arXiv:2504.08719}
  \item InfLLM-V2: Dense-Sparse Switchable Attention for Seamless Short-to-Long Adaptation (2025). \texttt{arXiv:2509.24663}
  \item Compact Recurrent Transformer with Persistent Memory (2025). \texttt{arXiv:2505.00929}
  \item StackTrans: From Large Language Model to Large Pushdown Automata Model (2025). \texttt{arXiv:2507.15343}
  \item PaPaformer: Language Model from Pre-trained Parallel Paths (2025). \texttt{arXiv:2508.00544}
  \item Gated Associative Memory: A Parallel O(N) Architecture for Efficient Sequence Modeling (2025). \texttt{arXiv:2509.00605}
  \item An extension of linear self-attention for in-context learning (2025). \texttt{arXiv:2503.23814}
  \item TPTT: Transforming Pretrained Transformers into Titans (2025). \texttt{arXiv:2506.17671}
  \item Inceptive Transformers: Enhancing Contextual Representations through Multi-Scale Feature Learning Across Domains and Languages (2025). \texttt{arXiv:2505.20496}
  \item Delayed Attention Training Improves Length Generalization in Transformer--RNN Hybrids (2025). \texttt{arXiv:2510.00258}
  \item ATLAS: Learning to Optimally Memorize the Context at Test Time (2025). \texttt{arXiv:2505.23735}
  \item From Fake Focus to Real Precision: Confusion-Driven Adversarial Attention Learning in Transformers (2025). \texttt{arXiv:2512.20661}
  \item Transformers as Multi-task Learners: Decoupling Features in Hidden Markov Models (2025). \texttt{arXiv:2506.01919}
  \item Self-Ablating Transformers: More Interpretability, Less Sparsity (2025). \texttt{arXiv:2505.00509}
  \item Hierarchical Resolution Transformers: A Wavelet-Inspired Architecture for Multi-Scale Language Understanding (2025). \texttt{arXiv:2509.20581}
  \item Short-Range Dependency Effects on Transformer Instability and a Decomposed Attention Solution (2025). \texttt{arXiv:2505.15548}
  \item Reservoir Computing as a Language Model (2025). \texttt{arXiv:2507.15779}
  \item Understanding the Skill Gap in Recurrent Language Models: The Role of the Gather-and-Aggregate Mechanism (2025). \texttt{arXiv:2504.18574}
  \item Beyond Parallelism: Synergistic Computational Graph Effects in Multi-Head Attention (2025). \texttt{arXiv:2507.02944}
  \item The underlying structures of self-attention: symmetry, directionality, and emergent dynamics in Transformer training (2025). \texttt{arXiv:2502.10927}
  \item Falcon-H1: A Family of Hybrid-Head Language Models Redefining Efficiency and Performance (2025). \texttt{arXiv:2507.22448}
\end{enumerate}

\end{document}
